% Figura 1 — Proporção de registros com MESPRENAT válido (códigos 1–9), por ano.
% Valores copiados do gráfico Plotly salvo em
% outputs/qualidade_geografica_2024/qualidade_sp_2014_2024.ipynb (seção 7).
% Denominador: todos os nascimentos do ano no recorte de residência materna.
\begin{tikzpicture}
\begin{axis}[
  estilo ipt,
  width=0.94\linewidth, height=6.1cm,
  xmin=2013.6, xmax=2024.4,
  ymin=93, ymax=100,
  xtick={2014,...,2024},
  xticklabel style={/pgf/number format/1000 sep={}},
  ytick={93,94,...,100},
  ymajorgrids,
  xlabel={Ano de nascimento},
  ylabel={MESPRENAT válido (\%)},
  legend style={at={(0.98,0.05)}, anchor=south east},
]
\addplot[serie principal, mark=*] coordinates {
  (2014,98.4480) (2015,98.4839) (2016,98.7980) (2017,99.0089) (2018,99.0324)
  (2019,98.9520) (2020,98.8038) (2021,98.9044) (2022,99.0966) (2023,99.1447)
  (2024,99.1547)};
\addlegendentry{Município de São Paulo (capital)}
\addplot[serie secundaria, mark=square*] coordinates {
  (2014,93.8857) (2015,93.8649) (2016,95.9476) (2017,97.2502) (2018,97.7259)
  (2019,97.9148) (2020,97.6013) (2021,97.9123) (2022,98.0692) (2023,98.3390)
  (2024,97.7061)};
\addlegendentry{Restante do estado de São Paulo}
% Rótulos seletivos: primeiro e último ano de cada série
\node[rotulo valor, anchor=south] at (axis cs:2014,98.4480) {98,45};
\node[rotulo valor, anchor=south] at (axis cs:2024,99.1547) {99,15};
\node[rotulo valor, anchor=north, yshift=-3pt] at (axis cs:2014,93.8857) {93,89};
\node[rotulo valor, anchor=north, yshift=-3pt] at (axis cs:2024,97.7061) {97,71};\end{axis}
\end{tikzpicture}
